\InterpolationArticle{7}
The case most frequent in practice is that in which \(a,b,c,d\mathinner{.\,.\,.}\) form an arithmetic progression. Setting the interval \(=1\), so that \(b=a+1,\ c=a+2\)
etc., the formula of Article 4 becomes
\[
\begin{aligned}
T={}&A\\
&+(t-1)(B-A)\\
&+(t-1)(t-2)\left(\frac{C-2B+A}{2}\right)\\
&+(t-1)(t-2)(t-3)\left(\frac{D-3C+3B-A}{6}\right)\\
&+\text{etc.}
\end{aligned}
\]
or
\[\begin{aligned}
T&=A+A'(t-1)+A''(t-1)(t-2)\\&+A'''(t-1)(t-2)(t-3)+\mathinner{.\,.}
\end{aligned}\]
where \(A',A'',A'''\) etc. are calculated by the following algorithm
\[
\begin{array}{lll}
A'=B-A,&B'=C-B,&C'=D-C\quad\text{etc.}\\
2A''=B'-A',&2B''=C'-B',&2C''=D'-C'\quad\text{etc.}\\
3A'''=B''-A'',&3B'''=C''-B'',&3C'''=D''-C''\quad\text{etc.}
\end{array}
\]
etc. This formula agrees completely with the usual formula for interpolation by differences.

\Needspace{42\baselineskip}
\InterpolationArticle{8}
If, for sufficiently many values of \(x\) proceeding in the arithmetic sequence
\(a,a+1,a+2\mathinner{.\,.\,.}\), the corresponding values of the function \(X\) are known, so that a sequence of \(m\)
successive values of \(x\) can be chosen at will for the calculation of \(T\):
the question arises, \emph{which} values it is best to prefer for this purpose, if we do not wish to use more than \(m\) or to go beyond the \(m-1^{\mathrm{th}}\) difference?
\SourceLeaf{286}
Clearly this must be decided so that the error arising from the difference of the \(m^{\mathrm{th}}\) order is
as small as possible; this error is
\[
=(t-a)(t-a-1)(t-a-2)\mathinner{.\,.\,.\,.}(t-a-m+1)\frac{\Delta}{1.2.3\mathinner{.\,.\,.} m}
\]
if the terms corresponding to \(x=a,a+1,a+2,\mathinner{.\,.\,.} a+m-1\) are used, or
\[
=(t-a-1)(t-a-2)(t-a-3)\mathinner{.\,.\,.\,.}(t-a-m)\frac{\Delta}{1.2.3\mathinner{.\,.\,.} m}
\]
if the terms corresponding to \(x=a+1,a+2,a+3,\mathinner{.\,.\,.} a+m\) are used, or
\[
=(t-a+1)(t-a)(t-a-1)\mathinner{.\,.\,.}(t-a-m+2)\frac{\Delta}{1.2.3\mathinner{.\,.\,.\,.}m}
\]
if the terms corresponding to \(x=a-1,a,a+1,\mathinner{.\,.\,.} a+m-2\) are used, denoting
by \(\Delta\) the difference of the \(m^{\mathrm{th}}\) order, which, since differences of higher orders
are disregarded, can here be considered constant. If, then, the \(m\) terms
corresponding to \(x=a,a+1,a+2,\mathinner{.\,.\,.} a+m-1\) are the most suitable, the first of those
three expressions must be the smallest, and thus, without regard to sign, \(t-a<\)
or at least not \(>t-a-m\), as well as \(t-a+1>\) or at least not \(<t-a-m+1\).
From the former condition it is easy to deduce that \(t-a-m\) must be a negative quantity lying between \(-\frac12m\) and \(-\infty\), whence \(t-a\) will lie between \(+\frac12m\) and \(-\infty\); from
the latter condition, however, \(t-a+1\) will be a positive quantity lying between \(\frac12m\) and
\(+\infty\); therefore \(t-a\) will lie between \(\frac12m-1\) and \(\frac12m\). It follows that, if \(m\) is an even number, the values to be used for interpolation must be chosen so that
the first half lie on one side of the desired term and the second half on the other; if, however,
\(m\) is odd, those terms must be chosen whose middle term lies as close as possible to the desired one.
Whenever, in this case, \(t\) is exactly midway between two consecutive values of \(x\), then, with respect to the error from the neglected difference of the \(m^{\mathrm{th}}\) order, generally speaking it makes no difference whether the term preceding or following the desired one is taken as the middle
of those to be used.

\InterpolationArticle{9}
\emph{Examples.} I. Suppose \(\log.\sin 24^{0}30'\) is to be found, when the logarithms of the sines are available for each degree, with fourth and higher differences neglected. In
this case, then, following the rules of the preceding article, we shall use the logarithms of the sines
of the arcs \(23^{0},24^{0},25^{0},26^{0}\). Denoting these by \(A,B,C,D\), the desired logarithm is obtained from the formula of Article 3 as
\SourceLeaf{287}
\[
=-\frac{1}{16}A+\frac{9}{16}B+\frac{9}{16}C-\frac{1}{16}D
\]
whose numerical value, from
\[
\begin{aligned}
A&=9{,}5918780\\
B&=9{,}6093133\\
C&=9{,}6259483\\
D&=9{,}6418420
\end{aligned}
\]
is found to be \(9{,}6177271.5\); the tables give \(9{,}6177270\).

If here, in place of \(\log.\sin 23^{0}\), we had used the logarithm of \(\sin 27^{0}=E=9{,}6570468\), the formula
\[
+\frac{5}{16}B+\frac{15}{16}C-\frac{5}{16}D+\frac{1}{16}E
\]
would have given \(9{,}6177267\,.\,375\).

For this case, the formula of the preceding article gives the error arising from the neglected fourth difference
\(\Delta\) as \(=\frac{3}{128}\Delta\) if \(A,B,C,D\) are used, but as \(=-\frac{5}{128}\Delta\) if \(B,C,D\)\raisebox{-0.3ex}{:}\,\(E\) are used; this error must be added to the calculated value. In our example,
\(\Delta=-0{,}0000066\); the correction calculated from this brings both calculations into agreement with the tables.

II. Suppose the time of the new moon on 16 June 1806 is to be calculated from the longitudes of the sun and moon for
each day, taking account of differences up to and including the fourth.
The Paris ephemerides give the following for local true noon

\begin{center}
\begin{tabular}{r@{\qquad}rrr@{\qquad}rrr}
&\multicolumn{3}{c}{Sun's longitude}&\multicolumn{3}{c}{Moon's longitude}\\
1806 June 14&\(82^{0}\)&\(39'\)&\(12''\)&\(53^{0}\)&\(19'\)&\(15''\)\\
15&83&36&30&67&29&43\\
16&84&33&48&81&59&52\\
17&85&31&7&96&44&9\\
18&86&28&24&111&35&30
\end{tabular}
\end{center}

\Needspace{15\baselineskip}
Here the given values of \(x\) are the differences between the longitudes of the sun and moon, and the corresponding values of \(X\) are the times, reckoned from noon on 16 June; the value of \(X\) corresponding to \(x=0\) is sought. Thus
\SourceLeaf{288}
\[
\begin{array}{r@{\;}c@{\;}r@{\;}r@{\;}r@{\quad}c@{\;}r@{\qquad}r@{\;}c@{\;}r}
a&=&-29^{0}&19'&57''&=&-105596''&A&=&-2\\
b&=&-16&6&47&=&-58007&B&=&-1\\
c&=&-2&33&56&=&-9236&C&=&0\\
d&=&+11&13&2&=&+40382&D&=&+1\\
e&=&+25&7&6&=&+90426&E&=&+2\\
t&=&0&&&&&&&
\end{array}
\]
\Needspace{23\baselineskip}
From this is calculated
\[
\begin{aligned}
\frac{(t-b)(t-c)(t-d)(t-e)}{(a-b)(a-c)(a-d)(a-e)}&=+0{,}0149084\\[5pt]
\frac{(t-a)(t-c)(t-d)(t-e)}{(b-a)(b-c)(b-d)(b-e)}&=-0{,}1050660\\[5pt]
\frac{(t-a)(t-b)(t-d)(t-e)}{(c-a)(c-b)(c-d)(c-e)}&=+0{,}9624567\\[5pt]
\frac{(t-a)(t-b)(t-c)(t-e)}{(d-a)(d-b)(d-c)(d-e)}&=+0{,}1434435\\[5pt]
\frac{(t-a)(t-b)(t-c)(t-d)}{(e-a)(e-b)(e-c)(e-d)}&=-0{,}0157429
\end{aligned}
\]
Multiplying these coefficients\HistoricalNote{23} respectively by \(-2,-1,0,1,2\), the sum of the products
is \(=+0{,}1872069\); hence the new moon will be on 16 June, at \(4^{\mathrm h}29'34''7\) Paris true
time, or \(4^{\mathrm h}29'41''4\) mean time; the Connaissance des tems gives \(4^{\mathrm h}27'42''\).

